% ==== Introduction ====
\section{Introduction}
\begin{frame}{Exception Handling: definition (1)}
  \begin{block}{Definition from Wikipedia, ``Exception handling'' article}
    \begin{center}
      ``{\ttfamily\small Exception handling is the process of responding to the occurrence of exceptions \textemdash{}anomalous or exceptional conditions requiring special processing\textemdash{} during the execution of a program.  In general, an exception breaks the normal flow of execution and executes a pre-registered exception handler.}''
    \end{center}
  \end{block}
  \vfill

  \begin{itemize}
    \item \textbf{In programming languages:} an exception object is prepared, which is then said to be \inlineCode{throw}n; control flow transfers to the inner-most \inlineCode{catch} block that can handle it
  \end{itemize}

  \begin{onlyenv}<2>
    \TZoverlayLayer{}{\fontsize{32}{32}\selectfont\smile}
  \end{onlyenv}
\end{frame}

\begin{frame}{Exception Handling: definition (2)}
  \begin{columns}
    \column{0.5\textwidth}
    Roughly, for C++\ldots{}
    \begin{enumerate}
    \item \textbf{User code:} Allocate memory for exception of type $T$ and throw it.
    \item \textbf{Runtime:} Identify the proper catch block ($Catch_T$).
    \item \textbf{Runtime:} Undo (unwind) affected stack frames; clean up as needed (dtors)
    \item \textbf{Runtime:} Transfer control to $Catch_T$.
    \end{enumerate}

    \column{0.5\textwidth}
    \begin{center}
      \resizebox{0.9\columnwidth}{!}{\tikz[
  start chain=c going below,node distance=0pt,
  >={Latex},
  annotation/.style={anchor=west,text width=48pt},
  annotation edge/.style={->,densely dotted},
  every node/.style={font={\footnotesize}},every label/.style={font={\footnotesize}}
]{
  \node(may_throw) [stack_frame={\#0}{black},on chain=c] {may\_throw()};
  \node(do_something_that_needs_cleanups) [stack_frame={\#1}{red},on chain=c] {do\_something\_that\_needs\_\\ cleanups()};
  \node(catches_and_rethrows) [stack_frame={\#2}{red},on chain=c] {catches\_and\_rethrows()};
  \node(main) [stack_frame={\#3}{red},on chain=c] {main()};

  \draw[annotation edge] (may_throw.east) .. controls +(15pt,-30pt) .. (catches_and_rethrows.east)
    node[annotation,very near start] {throws}
    node[annotation,very near end] {catches all};
  \draw[annotation edge] (catches_and_rethrows.east) .. controls +(15pt,-15pt) .. (main.east)
    node[annotation,very near start] {rethrows}
    node[annotation,very near end] {catches\\ {\ttfamily std::runtime\_error}};

  \node[below=10pt of main] {\textcolor{red!50!black!50}{\large$\bullet$}~needs cleanup};
}
}
    \end{center}
  \end{columns}

  \begin{onlyenv}<2>
    \overlayLayer{\begin{itemize}
      \item This is \alert{ABI}-dependent!
      \item Runtime needs additional information emitted by the compiler!
    \end{itemize}}
  \end{onlyenv}
\end{frame}

\begin{frame}{Let's talk about ABIs}
  \begin{block}{Informal definition}
    An ABI (Application Binary Interface) defines the low-level (machine-code-level) interfaces\\[1ex]
    \begin{tabular}{|p{0.47\textwidth}|p{0.47\textwidth}|}
      \hline
      \textbf{ABI}                                    & \textbf{API} \\
      \hline
      \hline
      low-level (binary) interface                    & high-level (source-code) interface \\
      \hline
      E.g., data layout, calling convention, \ldots{} & E.g., function name, parameter types, return type, \ldots{}\\
      \hline
    \end{tabular}
  \end{block}
  \vfill

  Well-known C++ ABIs:\begin{itemize}
    \item MSVC, used by Microsoft Visual C++
    \item Itanium C++ ABI, likely used by everything else (incl. \texttt{clang} / \texttt{gcc})\footnote{\url{https://itanium-cxx-abi.github.io/cxx-abi/}}
  \end{itemize}

  \begin{onlyenv}<2>
    \overlayLayer{\underline{A C++ ABI usually details}\begin{itemize}
      \item Exception handling
      \item Layout of RTTI structures
      \item Function call 'lowering' (e.g. \texttt{this} argument) --- \alert{but usually no calling convention!}
      \item Name mangling
    \end{itemize}}
  \end{onlyenv}
\end{frame}

\begin{frame}{Let's talk about ABIs: Calling Conventions}
  \begin{columns}
    \column{0.66\textwidth}
    A calling convention specifies\ldots{}
    \begin{itemize}
      \item \textbf{Function arguments / return values:} register / stack~(mem)?
      \item Which CPU registers are \textbf{scratch}, \textbf{caller-saved}, \textbf{callee-saved}?
      \item And, possibly, other stack frame layout details
        \bigskip
    \end{itemize}
    \vfill
    May be specified by a complementary ABI:
    \begin{itemize}
      \item \textbf{x86\_64, UNIX-based:} System V ABI\footnote{https://refspecs.linuxbase.org/}.
      \item \textbf{x86\_64, MSFT-based:} Microsoft x64 ABI
      \item \textbf{AArch64:} AAPCS64 (AArch64 Prodedure Call Standard)\footnote{https://github.com/ARM-software/abi-aa/blob/main/aapcs64/aapcs64.rst}.
    \end{itemize}

    \column{0.33\textwidth}
    \begin{center}
      \resizebox{0.95\columnwidth}{!}{\tikz[
  start chain=c going above,node distance=0pt,
  >={Latex},
  every node/.style={font={\footnotesize}}
]{
  \node(prev_LR) [stack_frame_element={20pt}{black},on chain=c,opacity=0.25] {(pushed) LR};
  \node(prev_FP) [stack_frame_element={20pt}{black},on chain=c,opacity=0.25] {(pushed) FP};
  %
  \draw[decorate,decoration=brace] ([xshift=-20pt]prev_LR.south west) -- ([xshift=-20pt,yshift=-1pt]prev_FP.north west)
    node[midway,anchor=east,text width=50pt,align=center] {Prev. stack frame};

  \node(LR) [stack_frame_element={20pt}{red},on chain=c] {(pushed) LR};
  \node(FP) [stack_frame_element={20pt}{red},on chain=c] {(pushed) FP};
  \node(X19) [stack_frame_element={20pt}{red},on chain=c] {(pushed) X19};
  \node(X20) [stack_frame_element={20pt}{red},on chain=c] {(pushed) X20};
  \node(X21) [stack_frame_element={20pt}{red},on chain=c] {(pushed) X21};
  \node(X22) [stack_frame_element={20pt}{red},on chain=c] {(pushed) X22};
  \node(16Bstackspace) [stack_frame_element={20pt}{black},on chain=c] {\ldots};
  %
  \draw[decorate,decoration=brace] ([xshift=-20pt]LR.south west) -- ([xshift=-20pt,yshift=1pt]16Bstackspace.north west)
    node[midway,anchor=east,text width=50pt,align=center] {Stack frame for {\ttfamily main()}};
  \draw[->,densely dotted,black!50] ([xshift=-5pt]LR.south west) -- ++(0pt,1in)
    node[midway,anchor=south,rotate=90] {decreasing addresses};

  \draw[<-,densely dotted] (LR.south east) -- ++(10pt,0pt) node[at end,anchor=west] {FP+0};
  \draw[<-,densely dotted] (FP.south east) -- ++(10pt,0pt) node[at end,anchor=west] {FP+8};
  \draw[<-,densely dotted] (X19.south east) -- ++(10pt,0pt) node[at end,anchor=west] {FP+16};
  \draw[<-,densely dotted] (X20.south east) -- ++(10pt,0pt) node[at end,anchor=west] {FP+24};
  \draw[<-,densely dotted] (X21.south east) -- ++(10pt,0pt) node[at end,anchor=west] {FP+32};
  \draw[<-,densely dotted] (X22.south east) -- ++(10pt,0pt) node[at end,anchor=west] {FP+40};
  %
  \draw[<-,densely dotted] (16Bstackspace.north east) .. controls ++(0pt,10pt) .. +(20pt,20pt) node[at end,anchor=west] {SP};
}
}
    \end{center}
  \end{columns}
\end{frame}
